\documentclass[11pt]{article}
\usepackage[margin=1in]{geometry}
\usepackage{times}
\usepackage{booktabs}
\usepackage{multirow}
\usepackage{graphicx}
\usepackage{amsmath}
\usepackage{hyperref}
\usepackage{natbib}

\title{Progress Report: Clarification-Seeking in Biomedical LLMs \\
Under Controlled PICO Ambiguity}

\author{
  Mihir \\
  % Add co-authors, affiliation, etc.
}

\date{March 2026}

\begin{document}
\maketitle

\begin{abstract}
We present preliminary results on a benchmark for testing whether large language models detect missing information in biomedical RCT abstracts and ask clarifying questions rather than guessing.
Using programmatic PICO slot deletion on the EBM-NLP dataset, we evaluate three models on missing-slot detection (Task~1) and demonstrate a full clarification pipeline with GPT-4o (Tasks~2--3).
Initial findings suggest that model scale and instruction-tuning are stronger predictors of detection ability than parameter count alone, and that a detect-then-clarify pipeline yields meaningful extraction improvements (+0.227 token-level F1).
These results are preliminary; key model comparisons (BioMistral-7B, Llama-3.1-8B) remain pending due to hardware constraints.
\end{abstract}

\section{Introduction}

Evidence-based medicine relies on structured extraction of PICO elements---Population (P), Intervention (I), and Outcome (O)---from randomized controlled trial (RCT) abstracts.
When abstracts are incomplete or ambiguous, practitioners and NLP systems must decide whether to guess or seek clarification.

We investigate whether LLMs can (1)~detect which PICO slot is missing from a manipulated abstract, (2)~generate a targeted clarifying question, and (3)~improve their extraction after receiving the missing information.
This framing connects information-seeking behavior in LLMs to a concrete downstream task with established evaluation infrastructure.

Our benchmark uses the EBM-NLP dataset \citep{nye2018corpus} with span-level \texttt{[REDACTED]} masking to create controlled ambiguity.
Unlike sentence-level deletion, span-level masking isolates the manipulation to a single PICO element without disrupting surrounding context.

\section{Experimental Setup}

\subsection{Dataset}

We use 189 test examples from EBM-NLP, with 75\% masked (one of P/I/O replaced with \texttt{[REDACTED]}) and 25\% unmasked controls.
Masking is span-level: only the annotated PICO span is replaced, preserving sentence structure.
We scope to three slots (P, I, O), excluding the Comparator slot due to noisy and inconsistent annotations in EBM-NLP.

\subsection{Tasks}

\begin{itemize}
  \item \textbf{Task~1 (Detection):} Given a (possibly masked) abstract, classify which PICO slot is missing from \{P, I, O, none\}.
  \item \textbf{Task~2 (Clarification):} For masked abstracts, generate a single targeted clarifying question about the missing element.
  \item \textbf{Task~3 (Pipeline):} Extract PICO elements before and after receiving the withheld information, measuring whether clarification improves extraction quality.
\end{itemize}

\subsection{Models}

For Task~1, we evaluate three models spanning different scales:
\begin{itemize}
  \item \textbf{GPT-4o}: Large proprietary model accessed via API.
  \item \textbf{Phi-3-mini-4k-instruct} (3.8B): Small instruction-tuned model, general-purpose.
  \item \textbf{SmolLM2-1.7B-Instruct} (1.7B): Smallest instruction-tuned baseline.
\end{itemize}

Each model is evaluated in zero-shot and few-shot (4 exemplars) settings.
Tasks~2 and~3 are currently evaluated with GPT-4o only.

We additionally compare against a teammate's fine-tuned BioBERT classifier (110M parameters), though this comparison is approximate: BioBERT uses 50/50 masking with sentence-level deletion, whereas our benchmark uses 75/25 masking with span-level deletion.

\paragraph{Hardware limitations.}
All local inference runs on Apple Silicon with 16\,GB unified memory.
BioMistral-7B (7B, $\sim$14\,GB in float16) exceeded the MPS buffer limit and could not be loaded.
Llama-3.1-8B-Instruct ($\sim$16\,GB in float16) also does not fit without quantization.
We plan to evaluate both models for the final report using either quantized inference or an external API.

\section{Preliminary Results}

\subsection{Task~1: Missing-Slot Detection}

Table~\ref{tab:task1} shows accuracy and macro-F1 across all evaluated models and settings.

\begin{table}[h]
\centering
\small
\begin{tabular}{llrcc}
\toprule
\textbf{Model} & \textbf{Setting} & \textbf{Params} & \textbf{Acc} & \textbf{M-F1} \\
\midrule
GPT-4o          & Zero-shot &       --- & 76.2 & 70.7 \\
GPT-4o          & Few-shot  &       --- & \textbf{83.1} & \textbf{82.0} \\
\midrule
Phi-3-mini      & Zero-shot &     3.8B & 48.7 & 42.1 \\
Phi-3-mini      & Few-shot  &     3.8B & 51.9 & 48.7 \\
\midrule
SmolLM2         & Zero-shot &     1.7B & 25.4 & 10.1 \\
SmolLM2         & Few-shot  &     1.7B & 29.6 & 24.7 \\
\midrule
BioBERT$^\dagger$  & Supervised &  110M & 82.5 & 78.0 \\
\bottomrule
\end{tabular}
\caption{Task~1 results. $^\dagger$BioBERT uses a different evaluation setup (50/50 masking ratio, sentence-level deletion) and is fine-tuned rather than prompted; included for reference only. Per-class metrics allow comparison despite different class distributions.}
\label{tab:task1}
\end{table}

\paragraph{Scaling trend.}
Performance increases monotonically with model scale: SmolLM2 (1.7B) performs near or below random, Phi-3-mini (3.8B) reaches roughly 50\% accuracy, and GPT-4o achieves over 80\% with few-shot prompting.
This ordering holds in both zero-shot and few-shot settings.

\paragraph{The none-class problem.}
Table~\ref{tab:none} reveals a systematic failure mode across prompted models: difficulty with the ``none'' class (unmasked abstracts where nothing is missing).

\begin{table}[h]
\centering
\small
\begin{tabular}{llccc}
\toprule
\textbf{Model} & \textbf{Setting} & \textbf{None-P} & \textbf{None-R} & \textbf{None-F1} \\
\midrule
GPT-4o     & Zero-shot & 1.00 & 0.17 & 0.29 \\
GPT-4o     & Few-shot  & 1.00 & 0.54 & 0.70 \\
Phi-3-mini & Zero-shot & 0.00 & 0.00 & 0.00 \\
Phi-3-mini & Few-shot  & 0.64 & 0.33 & 0.44 \\
SmolLM2    & Zero-shot & 0.25 & 1.00 & 0.41 \\
SmolLM2    & Few-shot  & 0.78 & 0.15 & 0.25 \\
\midrule
BioBERT$^\dagger$ & Supervised & --- & $\sim$0.91 & --- \\
\bottomrule
\end{tabular}
\caption{None-class (no missing slot) precision, recall, and F1. BioBERT recall estimated from confusion matrix. $^\dagger$Different setup.}
\label{tab:none}
\end{table}

Zero-shot GPT-4o and Phi-3-mini rarely predict ``none,'' defaulting to labeling some slot as missing even when none is.
SmolLM2 zero-shot shows the opposite pathology: it predicts ``none'' for everything (100\% recall, 25\% precision), consistent with random-level performance.
Few-shot examples partially correct this bias across all models, with GPT-4o reaching 0.70 none-F1.
By contrast, BioBERT's supervised training yields $\sim$91\% none-class recall, suggesting that the none-class problem is specific to prompted (rather than fine-tuned) approaches.

\subsection{Task~2: Clarification Question Quality}

We evaluate GPT-4o's clarifying questions using an LLM-as-judge protocol with three rubric criteria scored 1--5 (Table~\ref{tab:task2}).

\begin{table}[h]
\centering
\small
\begin{tabular}{lc}
\toprule
\textbf{Criterion} & \textbf{Mean Score} \\
\midrule
Slot targeting    & 4.75 / 5 \\
Specificity       & 4.36 / 5 \\
No assumptions    & 4.87 / 5 \\
\bottomrule
\end{tabular}
\caption{Task~2 LLM-as-judge evaluation (GPT-4o, $n = 141$).}
\label{tab:task2}
\end{table}

GPT-4o generates high-quality clarifying questions that target the correct missing slot (4.75), are specific enough for concise answers (4.36), and avoid introducing assumptions not present in the abstract (4.87).
These scores are from a single model judging itself, which is a limitation; we plan to add human evaluation for the final report.

\subsection{Task~3: Clarification Pipeline}

Table~\ref{tab:task3} shows the full detect--clarify--extract pipeline on GPT-4o.

\begin{table}[h]
\centering
\small
\begin{tabular}{lc}
\toprule
\textbf{Stage} & \textbf{Token-level PICO F1} \\
\midrule
Turn 1 (masked abstract)         & 0.298 \\
Turn 2 (after clarification)     & 0.524 \\
Oracle (full abstract)           & 0.478 \\
\midrule
\textbf{Delta (Turn 2 $-$ Turn 1)} & \textbf{+0.227} \\
\bottomrule
\end{tabular}
\caption{Task~3 PICO extraction F1 across pipeline stages (GPT-4o, $n = 141$). Macro-averaged over P, I, O slots.}
\label{tab:task3}
\end{table}

Providing the withheld information after a clarification turn yields a +0.227 improvement in token-level PICO F1.
Notably, turn~2 performance (0.524) slightly exceeds the oracle baseline (0.478), where the model extracts from the full unmasked abstract in a single turn.
This may reflect the benefit of the two-turn interaction: the model's attention is primed toward the previously missing slot, leading to more complete extraction.

\section{Discussion}

\paragraph{What these results suggest.}
The preliminary findings support three tentative conclusions: (1)~scale and instruction-tuning quality are the primary drivers of prompted missing-slot detection; (2)~the none-class represents a systematic weakness of prompted approaches that few-shot examples only partially mitigate; and (3)~a detect-then-clarify pipeline can meaningfully improve downstream extraction.

\paragraph{What remains uncertain.}
These results are preliminary and have notable gaps:

\begin{itemize}
  \item \textbf{BioMistral-7B} is absent due to hardware constraints. This model is critical for testing whether biomedical domain pretraining improves detection ability through prompting alone---the current results cannot distinguish ``scale matters'' from ``instruction-tuning quality matters,'' since GPT-4o has both.
  \item \textbf{Llama-3.1-8B-Instruct}, our originally proposed open-weight instruction-tuned model, also could not be evaluated locally. Running it via API or quantized inference is planned for the final report.
  \item \textbf{Task~2 judge bias}: GPT-4o judging its own outputs likely inflates scores. Human evaluation or a different judge model would provide a more reliable signal.
  \item \textbf{Single API model for Tasks~2--3}: We have not yet tested the full pipeline with open-weight models, which may show qualitatively different clarification behavior.
\end{itemize}

\section{Remaining Work}

For the final paper, we plan to:
\begin{enumerate}
  \item Evaluate BioMistral-7B and Llama-3.1-8B-Instruct using quantized local inference or an external API.
  \item Run Tasks~2 and~3 on at least one additional model beyond GPT-4o.
  \item Add human evaluation for Task~2 to complement the LLM-as-judge protocol.
  \item Analyze error patterns by slot type and abstract characteristics.
  \item Compare the masking-ratio sensitivity claim (75/25 vs.\ 50/50) if time permits.
\end{enumerate}

\bibliographystyle{plainnat}
\bibliography{references}

\end{document}
